% Generated draft: requires mathematical and source review.
\documentclass[11pt]{article}
\usepackage[margin=25mm]{geometry}
\usepackage{fontspec}
\setmainfont{Latin Modern Roman}
\usepackage{amsmath,amssymb,mathtools,unicode-math}
\setmathfont{Latin Modern Math}
\usepackage{enumitem,needspace,xurl,hyperref}
\setlength{\parindent}{0pt}
\setlength{\parskip}{5pt}
\setlength{\emergencystretch}{3em}
\newcommand{\N}{\mathbb{N}}
\newcommand{\Z}{\mathbb{Z}}
\newcommand{\Q}{\mathbb{Q}}
\newcommand{\R}{\mathbb{R}}
\newcommand{\C}{\mathbb{C}}
\newcommand{\F}{\mathbb{F}}
\newcommand{\A}{\mathbb{A}}
\newcommand{\PP}{\mathbb{P}}
\newcommand{\E}{\mathbb{E}}
\newcommand{\eps}{\varepsilon}
\newcommand{\dd}{\,\mathrm{d}}
\begin{document}
\section*{UnsolvedMath ID: 801001}
\textbf{Source problem code:} catalogue entry
\subsection{Definitions and mathematical statement}
Under the standard planar interpretation of the term, a \emph{two-point set} is a set \(X\subseteq\R^2\) meeting every affine line in exactly two distinct points. Here \(\R^2\) has its usual Euclidean topology, and a set is \emph{Borel} if it belongs to the smallest \(\sigma\)-algebra containing all open sets. An \(F_\sigma\) set is a countable union of closed sets.

Under this interpretation, the question is whether
\[
\exists X\subseteq\R^2\quad\left[X\text{ is Borel}\ \land\ \forall p\in\R^2\ \forall v\in\R^2\setminus\{0\},\quad \left|X\cap\{p+tv:t\in\R\}\right|=2\right].
\]
The supplied excerpt does not define ``two-point set'' or explicitly specify the ambient space. Thus the planar interpretation above is an explicit interpretation requiring source confirmation.

The supplied record states that every \(F_\sigma\) set is Borel and invokes Larman's theorem as ruling out the \(F_\sigma\) possibility while leaving the Borel question unsettled. This is preserved as supplied background attribution; the excerpt provides neither the theorem's precise hypotheses nor an original citation.
\subsection{Short English statement}
Does there exist a Borel subset of the Euclidean plane that meets every affine line in exactly two points? This uses the standard planar meaning of ``two-point set,'' which is not defined in the supplied excerpt. The registry labels the question open; that label is supplied source information, not an independently verified assessment.
\subsection{Sources}
\begin{itemize}
\item Local UnsolvedMath registry, permanent ID \(801001\), source problem code \texttt{catalogue entry}, title \emph{Open Problems in Mathematical Logic --- Problem 13}. Supplied upstream revision: \texttt{d371ee917ac9647fbb06cd6fc529deec6f3dda82}.
\item Barmpalias, Gao et al., \emph{Open Problems in Mathematical Logic} (2026), Problem 13, \url{https://arxiv.org/abs/2608.26628}. The registry records access on 2026-10-16 and extraction by \texttt{arxiv\_latex}; these are supplied provenance details, not a claim of independent source access. The registry attributes the problem to Barmpalias, Gao et al., with proposed year 2026.
\item The supplied statement attributes the exclusion of the \(F_\sigma\) case to Larman's theorem. No original reference or full theorem statement is supplied.
\end{itemize}
\end{document}
